% Generated draft: requires mathematical and source review.
\documentclass[11pt]{article}
\usepackage[margin=25mm]{geometry}
\usepackage{fontspec}
\setmainfont{Latin Modern Roman}
\usepackage{amsmath,amssymb,mathtools,unicode-math}
\setmathfont{Latin Modern Math}
\usepackage{enumitem,needspace,xurl,hyperref}
\setlength{\parindent}{0pt}
\setlength{\parskip}{5pt}
\setlength{\emergencystretch}{3em}
\newcommand{\N}{\mathbb{N}}
\newcommand{\Z}{\mathbb{Z}}
\newcommand{\Q}{\mathbb{Q}}
\newcommand{\R}{\mathbb{R}}
\newcommand{\C}{\mathbb{C}}
\newcommand{\F}{\mathbb{F}}
\newcommand{\A}{\mathbb{A}}
\newcommand{\PP}{\mathbb{P}}
\newcommand{\E}{\mathbb{E}}
\newcommand{\eps}{\varepsilon}
\newcommand{\dd}{\,\mathrm{d}}
\begin{document}
\section*{UnsolvedMath ID: 30006030}
\textbf{Source problem code:} catalogue entry
\subsection{Definitions and mathematical statement}
For a smooth connected compact complex curve \(C\) of genus \(g\in\Z_{\geq 0}\), a nonzero meromorphic \(k\)-differential, with \(k\in\Z_{>0}\), is a nonzero meromorphic section \(\eta\) of \(K_C^{\otimes k}\), where \(K_C\) is the canonical line bundle. If its zeros and poles are supported at distinct marked points \(p_1,\ldots,p_N\), write
\[
\operatorname{div}(\eta)=\sum_{i=1}^{N}m_i p_i,
\qquad m_i\in\Z,
\qquad \sum_{i=1}^{N}m_i=k(2g-2).
\]
Positive and negative orders describe zeros and poles, respectively. This supplies standard terminology on smooth complex curves; the source does not specify the intended ground field, compactification, or treatment of singular curves. Moreover, the source profile \(A\) cannot simply be identified with the order vector \((m_i)\): its displayed genus-zero example uses shifted orders.

\textbf{Problem.} For arbitrary genus \(g\), positive integer \(k\), and an admissible zero--pole profile \(A\) in the source's intended convention, formulate a natural enumerative problem involving \(k\)-differentials subject to prescribed period or residue conditions, and determine its relationship to the \(k\)-leaky double Hurwitz number \(H_g^k(A)\). The intended goal is an enumerative interpretation for \(k>0\), extending the genus-zero similarities reported in the source. No specific equality, normalization, or correspondence is supplied.

The source gives a genus-zero example with \(k=1\) and
\[
A=(d,-b_1,\ldots,-b_n),\qquad b_j\geq 0,
\qquad d=n-1+\sum_{j=1}^{n}b_j.
\]
Its displayed divisor formula has an indexing defect. If the listed \(b_j\) are paired in order with \(p_{j+1}\), the formula reads
\[
\operatorname{div}(\eta)
=(d-1)p_1-\sum_{j=1}^{n}(b_j+1)p_{j+1}.
\]
Under this explicitly conditional reindexing, the total order is \(-2\), as required for a meromorphic one-form in genus zero. The coefficients must be integers for this to be a divisor; the excerpt states nonnegativity of the \(b_j\) but does not explicitly state integrality. The source labels the associated locus by the flattened notation \(\texttt{DR00 (A)}\), whose intended typography and ambient moduli space are not specified.

\textbf{Unresolved formulation requirements.} A precise counting problem still requires the admissibility convention for \(A\), a definition and normalization of \(H_g^k(A)\), the period or residue constraints and their allowed values, and the moduli space, equivalence relation, and counting weights. In particular, scalar multiples of a differential must be accounted for. For \(k=1\), residues and periods have their usual meanings for meromorphic one-forms, with periods taken along appropriate cycles avoiding poles. For \(k>1\), the intended notions of periods and residues require additional specification. The supplied excerpt does not give the genus-zero residue conditions, the resulting count, or a general-genus conjectural formula. Consequently this is a request to formulate an enumerative interpretation, rather than a fully specified numerical identity.
\subsection{Short English statement}
The problem asks for a counting interpretation of \(H_g^k(A)\) using meromorphic \(k\)-differentials with prescribed periods or residues in arbitrary genus. The supplied source reports genus-zero similarities, but leaves the precise counting conditions, profile convention, and proposed relationship unspecified. Its partially solved label is supplied dated curation, not an independently verified mathematical conclusion.
\subsection{Sources}
\begin{itemize}
\item Local UnsolvedMath registry, permanent ID \(30006030\), source problem code \texttt{catalogue entry}, entitled \emph{k-Differential Interpretation of Leaky Double Hurwitz Numbers}. The supplied record lists proposed year 2024 and gives no named proposer.
\item Supplied report citation: \emph{Riemann Surfaces: Random, Flat, and Hyperbolic Geometry} (2025), Oberwolfach Reports, pp.\ 2369--2424, DOI \texttt{10.4171/owr/2024/41}. Report link, also listed among the supplied literature sources: \url{https://doi.org/10.4171/owr/2024/41}.
\item Supplied literature source labelled Cavalieri--Markwig--Schmitt: \url{https://doi.org/10.1017/fms.2025.26}. No article title or specific theorem is supplied in the record.
\item Supplied literature source labelled \emph{one-part leaky covers}: \url{https://arxiv.org/abs/2509.04335}. No authors or specific theorem are supplied in the record.
\item Literature assessment dated \textbf{2026-08-22}: the supplied curation labels the problem partially solved and reports that genus-zero residue-counting problems and structural comparisons with \(k\)-leaky numbers are known, while the arbitrary-genus enumerative problem and proposed relationship remain unproved. This is supplied dated curation derived from the OWR record, not independently verified here; its time-sensitive status and cited claims require independent checking.
\item The registry marks the statement \texttt{corrected\_verified}, explaining that only formula delimiters, notation, grammar, and typography were repaired. This is a supplied editorial designation and does not establish mathematical verification of this draft.
\end{itemize}
\end{document}
